% ====================================================================
% ch3v22_sec03_blend.tex — §3.3 혼합 음극 공통-μ 대정준 (R5 채택 = W1 기반 + W3 dQ/dV 합성식 그래프트)
% ====================================================================
\section{혼합 음극 --- 공통-$\mu$ 대정준 반전과 그 1차 근사 한계}\label{sec:si-blend}
흑연 다수에 실리콘 소수를 섞은 혼합(블렌드) 음극은 두 활물질이 \emph{같은 전위 손잡이를 공유}한다 ---
한 집전체$\cdot$한 전해질에 담긴 두 host 는 같은 전자 상($\phi_M$)과 같은 이온 상($\phi_S$)에 접하므로,
Part 0 의 전기화학 결선(\S\ref{sec:sm-electro} 식~\eqref{eq:sm-workbal})이 두 host 에 대해 같은 측정
전위 $V$ 로 닫히고, 곧 두 host 안 Li 의 화학퍼텐셜은 평형에서 \emph{하나의} $\mu$ 로 같아진다. 이것이
§3.1 이 앵커로 지목한 전하 보존의 전극 중립성(\S\ref{ssec:si-carry})의 혼합판이며, 본 절은 그
앵커 --- 다클래스 대정준 반전~\eqref{eq:sm-mc-balance}(\S\ref{sec:sm-mc}) --- 를 \emph{클래스곱에서
host 곱으로 한 줄 넓히는} 일이다. 재유도가 아니라 일반화다: 유일근을 보증한 요동 논증은 그대로 이월되고
(\S\ref{ssec:si-blend-derive}), $f_\mathrm{Si}\to0$ 에서 흑연 단독 결과가 그대로 회수된다
(\S\ref{ssec:si-blend-derive} 검산). 그러나 이 공통-$\mu$ 동시반응은 \emph{평형(OCV)에서 정확하되 유한
율속 블렌드 실측과는 정성적 편차}가 있으므로, \S\ref{ssec:si-blend-gs2} 가 그 편차를 정직 공백 GS-2 로
분리 진단한다.

\subsection{클래스곱에서 host 곱으로 --- 한 줄 일반화}\label{ssec:si-blend-derive}
\textbf{(a) 출발 --- 두 host, 한 저장조.} \S\ref{sec:sm-mc}(a)는 전극의 자리를 전이 $j$ 별 자리
클래스(site class)로 갈랐다: 클래스 $j$ 는 동등 자리 $M_j$ 개와 유효 자리 자유에너지
$\tilde\varepsilon_j(T)$ 를 갖고, 모든 클래스가 \emph{같은} 저장조와 입자를 교환하므로 화학퍼텐셜은
$\mu$ 하나뿐이었다. 혼합 음극은 자리 집합이 두 host 의 클래스 합집합일 뿐이다 --- host $\in\{\mathrm{gr},
\mathrm{Si}\}$ 마다 전이 $j=1,\dots,N_p^\mathrm{host}$ 의 클래스가 있고, 클래스 $j$ 는 자리 수
$M_j^\mathrm{host}$ 와 유효 자유에너지 $\tilde\varepsilon_j^\mathrm{host}(T)$ 를 갖는다. 결정적으로, 두
host 의 자리도 \emph{같은 하나의} 저장조와 교환한다(같은 전위 손잡이) --- 저장조 관점에서 흑연 클래스와
Si 클래스는 그저 $\mu$ 를 공유하는 독립 자리 클래스가 더 많아진 것이다. 가정 1(hard-core 배타)$\cdot$가정
2(자리 독립)는 host 사이로도 확장한다: host 안에서든 host 사이에서든 자리들은 독립이라 둔다(이 host 간
독립 가정의 물리 한계가 \S\ref{ssec:si-blend-gs2} 의 GS-2 다).

\textbf{(b) 연산 --- host 곱 인수분해.} 대정준 합은 독립 자리별 곱으로 갈라지고
(\S\ref{sec:sm-lattice} 와 같은 대수), 같은 클래스의 인수 $M_j^\mathrm{host}$ 개는 서로 같으므로,
\S\ref{sec:sm-mc}(b)의 클래스곱~\eqref{eq:sm-mc-factor} 이 host 라벨 하나를 더 달고 그대로 반복된다:
\begin{equation}
\Xi=\prod_{\mathrm{host}\in\{\mathrm{gr},\mathrm{Si}\}}\ \prod_{j=1}^{N_p^\mathrm{host}}
\big(\Xi_{1,j}^{\mathrm{host}}\big)^{M_j^\mathrm{host}},
\qquad
\ln\Xi=\sum_{\mathrm{host}}\sum_{j} M_j^\mathrm{host}\,\ln\Xi_{1,j}^\mathrm{host},
\qquad
\Xi_{1,j}^\mathrm{host}=1+e^{-\beta(\tilde\varepsilon_j^\mathrm{host}-\mu)}
\label{eq:blend-factor}
\end{equation}
--- 독립 부분계 분배함수의 표준 인수분해(\S\ref{sec:sm-mc} 가 문헌 근거와 함께 세운 그 대수)에서, 곱의
지표가 $j$ 하나에서 $(\mathrm{host},j)$ 쌍으로 늘어난 것뿐이다. 근사 경계도 \S\ref{sec:sm-mc}(b)의 두 겹을 이어받되 한 겹이
추가된다: (i) 클래스 \emph{안}의 $\Omega_j^\mathrm{host}$ 는 자기무모순 평균장으로만, (ii) 클래스
\emph{사이}(같은 host 내 staging 결합)의 상관은 빠지고, (iii) \emph{host 사이}의 상관 --- Si 팽창이
흑연 자리에 거는 기계 구속$\cdot$계면 결합 --- 도 이 곱에서 빠진다. 곧 식~\eqref{eq:blend-factor} 은 host
간 독립까지 포함한 평균장 수준에서 정확하다.

\textbf{(c) 중간식 --- 평균 입자수 $=$ host$\times$클래스 점유합.} 식~\eqref{eq:sm-gc} 의 미분
공식을 \eqref{eq:blend-factor} 의 로그에 적용하면 항별로 갈라져 \S\ref{sec:sm-mc}(c)의 계산이
$(\mathrm{host},j)$ 마다 반복된다:
\begin{equation}
\langle N\rangle=k_BT\,\frac{\partial\ln\Xi}{\partial\mu}
=\sum_{\mathrm{host}}\sum_{j} M_j^\mathrm{host}\,\theta_j^\mathrm{host}(\mu),
\qquad
\theta_j^\mathrm{host}=\frac{1}{1+e^{+\beta(\tilde\varepsilon_j^\mathrm{host}-\mu)}}
\label{eq:blend-occ}
\end{equation}
--- $\theta_j^\mathrm{host}$ 는 점유율~\eqref{eq:sm-mc-occ} 의 host 판이다. 저장된 Li 총량은 두 host 의
클래스 크기로 가중한 점유합이다.

\textbf{(d) 박스 --- 공통-$\mu$ 전하 보존 반전.} 입자수를 전하로 읽는다. 자리당 전하 $F/N_A$ 로 클래스
용량 $Q_j^\mathrm{host}\equiv(F/N_A)M_j^\mathrm{host}$, host 용량
$Q^\mathrm{host}\equiv\sum_jQ_j^\mathrm{host}$, 총용량 $Q\equiv\sum_\mathrm{host}Q^\mathrm{host}$ 를
두고, 추출(탈리튬화) 전하 분율을 $\bar x\equiv1-\langle N\rangle/\sum_{\mathrm{host},j}M_j^\mathrm{host}$
로 두면, \eqref{eq:blend-occ} 에서 진행률 $\xi_{\eq,j}^\mathrm{host}=1-\theta_j^\mathrm{host}$ 가
$\mu\leftrightarrow V$ 결선(\S\ref{sec:sm-electro})을 거쳐 host 별 평형 진행률
$\xi_{\eq,j}^\mathrm{host}(V,T)$~\eqref{eq:sm-logistic} 이 되고, 고정된 $\bar x$ 에서 이 등식을
만족시키는 \emph{하나의} 전위가 관측 평형 전위 $U_\mathrm{oc}(\bar x,T)$ 다:
\begin{equation}
\boxed{\;\sum_{\mathrm{host}\in\{\mathrm{gr},\mathrm{Si}\}}\ \sum_{j=1}^{N_p^\mathrm{host}}
Q_j^\mathrm{host}\,\xi_{\eq,j}^\mathrm{host}\bigl(U_\mathrm{oc},T\bigr)=Q\,\bar x,
\qquad
Q=Q_\mathrm{gr}+Q_\mathrm{Si},\quad f_\mathrm{Si}=\frac{Q_\mathrm{Si}}{Q}\;.}
\label{eq:blend-balance}
\end{equation}
이 식은 다클래스 반전~\eqref{eq:sm-mc-balance} 에서 단일 합 $\sum_j$ 를 이중 합
$\sum_\mathrm{host}\sum_j$ 로 바꾼 \emph{한 줄}이다 --- 두 식을 나란히 두면 바뀐 것은 지표의 범위뿐이고,
전하 회계$\cdot$결선$\cdot$반전의 논리는 글자까지 같다.\footnote{범위 좌표 규약 --- 선언$\cdot$스윕 범위는 업계 관행(공장의 질량 배합)에 맞춘 \emph{질량} 분율
$m_\mathrm{Si}\in[0,0.30]$(0--30 wt\%) 기준이다. 전하 보존식 내부 변수는 \emph{용량} 분율
$f_\mathrm{Si}=Q_\mathrm{Si}/Q$ 이며 둘은 비용량 환산
$f_\mathrm{Si}=m_\mathrm{Si}\,q_\mathrm{Si}\,/\,[\,m_\mathrm{Si}\,q_\mathrm{Si}+(1-m_\mathrm{Si})\,q_\mathrm{gr}\,]$($q$ $=$ 비용량)로
잇는다 --- $q_\mathrm{Si}\gg q_\mathrm{gr}$ 라 wt\% 수치보다 훨씬 크게 앉아, $[0,0.30]$ wt\% 는 용량 분율로
대체로 $f_\mathrm{Si}\approx0$--$0.8$ 에 대응한다(케이스별 $q_\mathrm{Si}$ 에 따라 30 wt\% 에서 $f_\mathrm{Si}{\approx}0.54$[원소 Si]--$0.78$[Si--C]; \S\ref{ssec:si-sic}$\cdot$\S\ref{ssec:si-blend-gs2} 의
10--30 wt\% 실측 앵커 포함). wt\% 기반 실측과의 대조는 이 환산을 거친다.} 반전의 지위도 그대로다: $\bar x$ 를 고정하고
전위를 푸는 고정-함량 반전이라 $U_\mathrm{oc}$ 가 음함수로 앉는 것은 논리 공백이 아니라 \emph{정의상
implicit formulation}(대정준$\to$정준 Legendre-켤레 반전)이며, 두 host 각각 전이가 여럿이라 항상
$N_p^\mathrm{gr}+N_p^\mathrm{Si}\ge2$ 겹침이므로 닫힌 역이 없어 \emph{수치 유일근}으로 푼다. $f_\mathrm{Si}$
는 그 반전에 들어가는 용량 배분 하나를 조절하는 손잡이이며(§3.5 코드의 토글), 두 host 의 서로 다른 전이
집합($U_j^\mathrm{host}\cdot w_j^\mathrm{host}$, §3.2)이 같은 $U_\mathrm{oc}$ 인자를 공유해 합산된다.
\begin{verifybox}
검산 세 건 --- (i) \emph{유일근(요동 양성) 이월.} 식~\eqref{eq:blend-occ} 을 $\mu$ 로 미분하면
$(\mathrm{host},j)$ 마다 요동--응답이 반복되어
$\partial\langle N\rangle/\partial\mu=\beta\sum_\mathrm{host}\sum_jM_j^\mathrm{host}\theta_j^\mathrm{host}
(1-\theta_j^\mathrm{host})=\beta\,\mathrm{var}(N)>0$ 이다 --- 독립 자리 분산의 가법이 host 사이로도 성립
(두 host 의 자리는 서로 독립)하므로, 흑연 host 에 Si host 를 더하는 일은 좌변에 \emph{양의 분산 항을 더
얹는} 것일 뿐이다. 곧 \S\ref{sec:sm-mc} 의 요동 양성 유일근 증명(식~\eqref{eq:sm-mc-fluc})이 재유도
없이 그대로 이월되어, 결선~\eqref{eq:sm-eqcond} 하에서 \eqref{eq:blend-balance} 좌변은 $U_\mathrm{oc}$
에 연속 순증이고 $U_\mathrm{oc}\to\mp\infty$ 에서 $0/Q$ 로 가 $\bar x\in(0,1)$ 마다 근이 유일하다(가드도
이월 --- 평균장 이중웰 $\Omega_j^\mathrm{host}>2RT$ 의 비단조 가지는 이 보증 밖이고, 실제 반전에 쓰는
전이별 logistic 은 항마다 강단조라 보증 안). (ii) \emph{$f_\mathrm{Si}\to0$ 회수(bit-exact).}
$f_\mathrm{Si}=Q_\mathrm{Si}/Q\to0$ 은 $Q_\mathrm{Si}\to0$ 이고, Si 클래스 용량이
$Q_j^\mathrm{Si}=Q_\mathrm{Si}\cdot(Q_j^\mathrm{Si}/Q_\mathrm{Si})\to0$ 이라 Si 이중합 항이 통째로 사라져
\begin{equation}
\eqref{eq:blend-balance}\ \xrightarrow{\,f_\mathrm{Si}\to0\,}\
\sum_{j=1}^{N_p^\mathrm{gr}}Q_j^\mathrm{gr}\,\xi_{\eq,j}^\mathrm{gr}(U_\mathrm{oc},T)=Q_\mathrm{gr}\,\bar x
\label{eq:blend-limit}
\end{equation}
--- $Q=Q_\mathrm{gr}$ 의 흑연 단독 다클래스 반전~\eqref{eq:sm-mc-balance} 로 \emph{글자까지} 되돌아간다.
이는 \S\ref{sec:sm-mc} 검산 (ii)의 $N_p{=}1$ 회수와 같은 정신의 확장 off 이며, §3.5 코드의
``$f_\mathrm{Si}{=}0$ $\to$ 흑연 단독 bit-exact'' 게이트의 \emph{이론적 짝}이다(수치 재현이 아니라 식이
같아짐). (iii) \emph{부호 읽기.} 결선~\eqref{eq:sm-eqcond}($s{=}+1$)로 $V\!\uparrow\Rightarrow$ 모든
$\theta_j^\mathrm{host}\!\downarrow\Rightarrow\sum_{\mathrm{host},j}Q_j^\mathrm{host}
\xi_{\eq,j}^\mathrm{host}\!\uparrow$ --- 두 host 어디서든 전위를 올리면 추출이 진행하고, $\bar x$ 고정에서
$\mu$ 를 푸는 일이 곧 $V$ 를 푸는 일이라 그 근이 공통 $U_\mathrm{oc}(\bar x,T)$ 다.
\end{verifybox}

관측 미분용량은 이 반전의 평형 종 합성으로 곧장 적힌다 --- 합산식~\eqref{eq:sum} 의 host 판이다:
\begin{equation}
\frac{\dd Q}{\dd V}\bigg|_{U_\mathrm{oc}}=C_\bg+\sum_{\mathrm{host}\in\{\mathrm{gr},\mathrm{Si}\}}\ \sum_{j}
\frac{Q_j^{\mathrm{host}}\,\xi_{\eq,j}^{\mathrm{host}}\bigl(1-\xi_{\eq,j}^{\mathrm{host}}\bigr)}{w_j^{\mathrm{host}}}
\label{eq:blend-dqdv}
\end{equation}
--- 두 host 의 종이 \emph{같은 전위 축} 위에 겹쳐 놓이고, 겹치는 전위 창에서 두 봉우리가 포갠다. 저-Si
블렌드 dQ/dV 에서 Si$\cdot$흑연 피크가 부분 분리로 관측되는 것(\S\ref{ssec:si-sic}, Si $10$
wt\%\cite{naboka_sic2021})이 이 겹침의 실측 접점이며, 배경 $C_\bg$ 는 전극 단위로 한 번만 더한다(host 별
중복 가산 금지).

\begin{keybox}
\textbf{공통-$\mu$ 앵커.} 혼합 음극의 중심식은 다클래스 반전~\eqref{eq:sm-mc-balance} 의 지표를
$\sum_j\to\sum_\mathrm{host}\sum_j$ 로 넓힌 \eqref{eq:blend-balance} 하나다 --- 요동 양성 유일근은
이월, $f_\mathrm{Si}\to0$ 흑연 회수는 식~\eqref{eq:blend-limit} 로 닫힌다. 관측 dQ/dV 는 두 host 전이 기여의 배경 위 선형 합~\eqref{eq:blend-dqdv}(\S\ref{sec:sum} 식~\eqref{eq:sum}
의 host 판)으로 얻으며, 이것이 §3.5 코드 합성의 규칙이다.
\end{keybox}

\begin{figure}[t]
\centering
\begin{tikzpicture}
% ============================================================
% (a) 전체 창 — 총 평형 dQ/dV 를 총용량 Q 로 정규화한 가족 (4곡선)
%     x = V_n [V] (0.03–0.62) , y = (dQ/dV)/Q [1/V] (면적=1)
% ============================================================
\begin{scope}[x=9.5cm,y=0.58cm]
% ---- 축 ----
\draw[-{Latex[length=1.6mm]}] (0.026,0) -- (0.026,7.95)
  node[above,font=\scriptsize] {$(\dd Q/\dd V)/Q$ [1/V]};
\draw[-{Latex[length=1.6mm]}] (0.026,0) -- (0.642,0)
  node[below,font=\scriptsize] {$V_n$ [V]};
\foreach \xx in {0.1,0.2,0.3,0.4,0.5}
  {\draw (\xx,0.09) -- (\xx,-0.09) node[below,font=\scriptsize] {\xx};}
\foreach \yy in {2,4,6}
  {\draw (0.0235,\yy) -- (0.0285,\yy) node[left,font=\scriptsize] {\yy};}
% ---- 가족 곡선 (흑백; 선종+명도로 구분) ----
% 0 wt%  f_Si=0  (흑연 단독, 기준선) — very thick 실선
\draw[very thick] plot[smooth] coordinates {(0.0300,2.2213) (0.0360,2.6753) (0.0420,3.1886) (0.0480,3.7540) (0.0540,4.3580) (0.0600,4.9795) (0.0660,5.5915) (0.0720,6.1625) (0.0780,6.6611) (0.0840,7.0604) (0.0900,7.3411) (0.0960,7.4944) (0.1020,7.5208) (0.1080,7.4279) (0.1140,7.2282) (0.1200,6.9362) (0.1260,6.5677) (0.1320,6.1396) (0.1380,5.6701) (0.1440,5.1785) (0.1500,4.6844) (0.1560,4.2061) (0.1620,3.7593) (0.1680,3.3554) (0.1740,3.0012) (0.1800,2.6980) (0.1860,2.4423) (0.1920,2.2268) (0.1980,2.0416) (0.2040,1.8765) (0.2100,1.7222) (0.2160,1.5717) (0.2220,1.4213) (0.2280,1.2704) (0.2340,1.1206) (0.2400,0.9750) (0.2460,0.8371) (0.2520,0.7097) (0.2580,0.5950) (0.2640,0.4940) (0.2700,0.4066) (0.2760,0.3322) (0.2820,0.2699) (0.2880,0.2181) (0.2940,0.1756) (0.3000,0.1409) (0.3160,0.0774) (0.3320,0.0421) (0.3480,0.0227) (0.3640,0.0122) (0.3800,0.0066) (0.3960,0.0035) (0.4120,0.0019) (0.4280,0.0010) (0.4440,0.0005) (0.4600,0.0003) (0.4760,0.0002) (0.4920,0.0001) (0.6200,0.0000)};
% 10 wt%  f_Si=0.230 — thick densely dashed
\draw[thick,densely dashed] plot[smooth] coordinates {(0.0300,1.7358) (0.0360,2.0880) (0.0420,2.4861) (0.0480,2.9247) (0.0540,3.3932) (0.0600,3.8757) (0.0660,4.3510) (0.0720,4.7953) (0.0780,5.1844) (0.0840,5.4973) (0.0900,5.7196) (0.0960,5.8442) (0.1020,5.8717) (0.1080,5.8081) (0.1140,5.6629) (0.1200,5.4473) (0.1260,5.1737) (0.1320,4.8550) (0.1380,4.5053) (0.1440,4.1396) (0.1500,3.7728) (0.1560,3.4192) (0.1620,3.0909) (0.1680,2.7967) (0.1740,2.5418) (0.1800,2.3272) (0.1860,2.1502) (0.1920,2.0050) (0.1980,1.8842) (0.2040,1.7796) (0.2100,1.6840) (0.2160,1.5919) (0.2220,1.5004) (0.2280,1.4085) (0.2340,1.3176) (0.2400,1.2297) (0.2460,1.1473) (0.2520,1.0723) (0.2580,1.0060) (0.2640,0.9490) (0.2700,0.9012) (0.2760,0.8618) (0.2820,0.8297) (0.2880,0.8040) (0.2940,0.7833) (0.3000,0.7666) (0.3160,0.7349) (0.3320,0.7134) (0.3480,0.6968) (0.3640,0.6845) (0.3800,0.6767) (0.3960,0.6724) (0.4120,0.6677) (0.4280,0.6570) (0.4440,0.6341) (0.4600,0.5954) (0.4760,0.5407) (0.4920,0.4741) (0.5080,0.4019) (0.5240,0.3304) (0.5400,0.2646) (0.5560,0.2074) (0.5720,0.1599) (0.5880,0.1216) (0.6040,0.0916) (0.6200,0.0685)};
% 20 wt%  f_Si=0.402 — thick densely dashdotted, black!75
\draw[thick,densely dashdotted,black!75] plot[smooth] coordinates {(0.0300,1.3729) (0.0360,1.6490) (0.0420,1.9610) (0.0480,2.3047) (0.0540,2.6720) (0.0600,3.0504) (0.0660,3.4237) (0.0720,3.7733) (0.0780,4.0804) (0.0840,4.3288) (0.0900,4.5073) (0.0960,4.6105) (0.1020,4.6389) (0.1080,4.5971) (0.1140,4.4927) (0.1200,4.3343) (0.1260,4.1315) (0.1320,3.8946) (0.1380,3.6345) (0.1440,3.3628) (0.1500,3.0913) (0.1560,2.8310) (0.1620,2.5912) (0.1680,2.3791) (0.1740,2.1984) (0.1800,2.0500) (0.1860,1.9318) (0.1920,1.8393) (0.1980,1.7665) (0.2040,1.7072) (0.2100,1.6555) (0.2160,1.6071) (0.2220,1.5594) (0.2280,1.5118) (0.2340,1.4649) (0.2400,1.4202) (0.2460,1.3792) (0.2520,1.3433) (0.2580,1.3132) (0.2640,1.2893) (0.2700,1.2710) (0.2760,1.2576) (0.2820,1.2483) (0.2880,1.2419) (0.2940,1.2376) (0.3000,1.2343) (0.3160,1.2265) (0.3320,1.2153) (0.3480,1.2008) (0.3640,1.1871) (0.3800,1.1777) (0.3960,1.1724) (0.4120,1.1654) (0.4280,1.1473) (0.4440,1.1078) (0.4600,1.0402) (0.4760,0.9448) (0.4920,0.8285) (0.5080,0.7023) (0.5240,0.5774) (0.5400,0.4624) (0.5560,0.3625) (0.5720,0.2794) (0.5880,0.2125) (0.6040,0.1601) (0.6200,0.1197)};
% 30 wt%  f_Si=0.535 — thick densely dotted, black!58
\draw[thick,densely dotted,black!58] plot[smooth] coordinates {(0.0300,1.0913) (0.0360,1.3084) (0.0420,1.5536) (0.0480,1.8237) (0.0540,2.1124) (0.0600,2.4101) (0.0660,2.7043) (0.0720,2.9803) (0.0780,3.2238) (0.0840,3.4222) (0.0900,3.5667) (0.0960,3.6533) (0.1020,3.6824) (0.1080,3.6576) (0.1140,3.5847) (0.1200,3.4707) (0.1260,3.3230) (0.1320,3.1495) (0.1380,2.9589) (0.1440,2.7602) (0.1500,2.5626) (0.1560,2.3746) (0.1620,2.2036) (0.1680,2.0550) (0.1740,1.9319) (0.1800,1.8349) (0.1860,1.7623) (0.1920,1.7107) (0.1980,1.6752) (0.2040,1.6510) (0.2100,1.6333) (0.2160,1.6188) (0.2220,1.6053) (0.2280,1.5920) (0.2340,1.5792) (0.2400,1.5679) (0.2460,1.5591) (0.2520,1.5535) (0.2580,1.5516) (0.2640,1.5532) (0.2700,1.5579) (0.2760,1.5647) (0.2820,1.5730) (0.2880,1.5817) (0.2940,1.5900) (0.3000,1.5973) (0.3160,1.6079) (0.3320,1.6046) (0.3480,1.5918) (0.3640,1.5770) (0.3800,1.5664) (0.3960,1.5603) (0.4120,1.5516) (0.4280,1.5278) (0.4440,1.4753) (0.4600,1.3854) (0.4760,1.2583) (0.4920,1.1035) (0.5080,0.9354) (0.5240,0.7690) (0.5400,0.6159) (0.5560,0.4828) (0.5720,0.3721) (0.5880,0.2831) (0.6040,0.2132) (0.6200,0.1594)};
% ---- 추세 주석: 흑연 축소(아래) · Si 성장(위) ----
\draw[-{Latex[length=1.4mm]},thick] (0.132,6.35) -- (0.132,3.55);
\node[font=\tiny,anchor=west,align=left] at (0.146,5.75) {흑연 staging 겹침\\envelope 축소 ($m_\mathrm{Si}\!\uparrow$)};
\draw[-{Latex[length=1.4mm]},thick] (0.320,0.25) -- (0.320,1.48);
\node[font=\tiny,anchor=west,align=left] at (0.345,1.92) {Si 기여 성장\\($m_\mathrm{Si}\!\uparrow$)};
% ---- 범례 (우상단 빈 영역) ----
\draw[very thick] (0.400,7.35) -- (0.437,7.35);
\node[font=\tiny,anchor=west] at (0.443,7.35) {$0$ wt\% ($f_\mathrm{Si}{=}0$, 흑연 단독)};
\draw[thick,densely dashed] (0.400,6.65) -- (0.437,6.65);
\node[font=\tiny,anchor=west] at (0.443,6.65) {$10$ wt\% ($f_\mathrm{Si}{=}0.23$)};
\draw[thick,densely dashdotted,black!75] (0.400,5.95) -- (0.437,5.95);
\node[font=\tiny,anchor=west] at (0.443,5.95) {$20$ wt\% ($f_\mathrm{Si}{=}0.40$)};
\draw[thick,densely dotted,black!58] (0.400,5.25) -- (0.437,5.25);
\node[font=\tiny,anchor=west] at (0.443,5.25) {$30$ wt\% ($f_\mathrm{Si}{=}0.54$)};
\node[font=\scriptsize] at (0.33,-1.05) {(a) 전체 창 --- 총 dQ/dV$/Q$};
\end{scope}
% ============================================================
% (b) Si 기여 확대 — host_contributions 의 Si 몫만(eq:blend-dqdv Si 이중합)/Q
%     x = V_n [V] (0.15–0.62) , y = (dQ/dV)_Si/Q [1/V]
% ============================================================
\begin{scope}[xshift=6.4cm,x=9.5cm,y=2.6cm]
\draw[-{Latex[length=1.6mm]}] (0.146,0) -- (0.146,1.74)
  node[above,font=\scriptsize] {$(\dd Q/\dd V)_\mathrm{Si}/Q$ [1/V]};
\draw[-{Latex[length=1.6mm]}] (0.146,0) -- (0.644,0)
  node[below,font=\scriptsize] {$V_n$ [V]};
\foreach \xx in {0.2,0.3,0.4,0.5}
  {\draw (\xx,0.02) -- (\xx,-0.02) node[below,font=\scriptsize] {\xx};}
\foreach \yy in {0.5,1.0,1.5}
  {\draw (0.1445,\yy) -- (0.1475,\yy) node[left,font=\scriptsize] {\yy};}
% 10 wt%
\draw[thick,densely dashed] plot[smooth] coordinates {(0.1500,0.1658) (0.1590,0.1882) (0.1680,0.2130) (0.1770,0.2402) (0.1860,0.2696) (0.1950,0.3012) (0.2040,0.3347) (0.2130,0.3698) (0.2220,0.4059) (0.2310,0.4426) (0.2400,0.4790) (0.2490,0.5143) (0.2580,0.5478) (0.2670,0.5786) (0.2760,0.6059) (0.2850,0.6292) (0.2940,0.6481) (0.3030,0.6624) (0.3120,0.6723) (0.3210,0.6782) (0.3300,0.6808) (0.3390,0.6809) (0.3480,0.6793) (0.3570,0.6770) (0.3660,0.6746) (0.3750,0.6725) (0.3840,0.6711) (0.3930,0.6700) (0.4020,0.6688) (0.4110,0.6666) (0.4200,0.6624) (0.4290,0.6552) (0.4380,0.6439) (0.4470,0.6278) (0.4560,0.6063) (0.4650,0.5797) (0.4740,0.5482) (0.4830,0.5126) (0.4920,0.4741) (0.5010,0.4337) (0.5100,0.3928) (0.5190,0.3523) (0.5280,0.3132) (0.5370,0.2763) (0.5460,0.2421) (0.5550,0.2107) (0.5640,0.1824) (0.5730,0.1572) (0.5820,0.1349) (0.5910,0.1154) (0.6000,0.0984) (0.6090,0.0837) (0.6180,0.0711)};
% 20 wt%
\draw[thick,densely dashdotted,black!75] plot[smooth] coordinates {(0.1500,0.2897) (0.1590,0.3290) (0.1680,0.3723) (0.1770,0.4197) (0.1860,0.4711) (0.1950,0.5263) (0.2040,0.5849) (0.2130,0.6462) (0.2220,0.7094) (0.2310,0.7734) (0.2400,0.8370) (0.2490,0.8988) (0.2580,0.9574) (0.2670,1.0111) (0.2760,1.0589) (0.2850,1.0996) (0.2940,1.1325) (0.3030,1.1575) (0.3120,1.1748) (0.3210,1.1852) (0.3300,1.1897) (0.3390,1.1899) (0.3480,1.1872) (0.3570,1.1831) (0.3660,1.1789) (0.3750,1.1753) (0.3840,1.1728) (0.3930,1.1709) (0.4020,1.1687) (0.4110,1.1649) (0.4200,1.1576) (0.4290,1.1450) (0.4380,1.1253) (0.4470,1.0971) (0.4560,1.0596) (0.4650,1.0130) (0.4740,0.9580) (0.4830,0.8958) (0.4920,0.8285) (0.5010,0.7579) (0.5100,0.6864) (0.5190,0.6157) (0.5280,0.5474) (0.5370,0.4829) (0.5460,0.4230) (0.5550,0.3683) (0.5640,0.3188) (0.5730,0.2748) (0.5820,0.2358) (0.5910,0.2017) (0.6000,0.1720) (0.6090,0.1463) (0.6180,0.1242)};
% 30 wt%
\draw[thick,densely dotted,black!58] plot[smooth] coordinates {(0.1500,0.3859) (0.1590,0.4381) (0.1680,0.4958) (0.1770,0.5590) (0.1860,0.6275) (0.1950,0.7010) (0.2040,0.7790) (0.2130,0.8607) (0.2220,0.9448) (0.2310,1.0301) (0.2400,1.1149) (0.2490,1.1972) (0.2580,1.2751) (0.2670,1.3468) (0.2760,1.4104) (0.2850,1.4646) (0.2940,1.5084) (0.3030,1.5417) (0.3120,1.5648) (0.3210,1.5786) (0.3300,1.5846) (0.3390,1.5848) (0.3480,1.5812) (0.3570,1.5758) (0.3660,1.5701) (0.3750,1.5654) (0.3840,1.5620) (0.3930,1.5595) (0.4020,1.5566) (0.4110,1.5516) (0.4200,1.5419) (0.4290,1.5250) (0.4380,1.4988) (0.4470,1.4612) (0.4560,1.4113) (0.4650,1.3493) (0.4740,1.2759) (0.4830,1.1932) (0.4920,1.1034) (0.5010,1.0095) (0.5100,0.9142) (0.5190,0.8200) (0.5280,0.7291) (0.5370,0.6432) (0.5460,0.5634) (0.5550,0.4905) (0.5640,0.4247) (0.5730,0.3659) (0.5820,0.3141) (0.5910,0.2686) (0.6000,0.2291) (0.6090,0.1948) (0.6180,0.1654)};
% 성장 방향 주석 + Si 창 안내
\draw[-{Latex[length=1.4mm]},thick] (0.240,0.55) -- (0.240,1.42);
\node[font=\tiny,anchor=west] at (0.152,1.62) {Si host 기여항 (eq~\eqref{eq:blend-dqdv} Si 몫); $0$ wt\%: 기여 없음};
\node[font=\tiny,anchor=west] at (0.255,1.30) {$m_\mathrm{Si}\!\uparrow$};
\node[font=\scriptsize] at (0.39,-0.235) {(b) Si 기여 확대};
\end{scope}
\end{tikzpicture}
\caption{블렌드 $f_\mathrm{Si}$(wt\% 스윕) 가족 평형 dQ/dV --- 좌표는 원소 Si(elemental) host 블렌드의
평형식~\eqref{eq:blend-dqdv} 실계산이다:
그 $|I|\!\to\!0$ 종
$\sum_{\mathrm{host}}\sum_j Q_j^\mathrm{host}\xi_{\eq,j}^\mathrm{host}(1-\xi_{\eq,j}^\mathrm{host})/w_j^\mathrm{host}$
을 $T{=}298.15$~K$\cdot$$C_\bg{=}0$ 에서 평가하고 총용량 $Q{=}Q_\mathrm{gr}{+}Q_\mathrm{Si}$ 로
정규화했다(곡선 아래 면적 $=1$ 규약, 곡선당 $\sim$$53$--$66$ 표본점). 질량$\leftrightarrow$용량 환산은
§\ref{ssec:si-blend-derive} 각주의 환산 규약
$f_\mathrm{Si}=m_\mathrm{Si}q_\mathrm{Si}/[m_\mathrm{Si}q_\mathrm{Si}+(1{-}m_\mathrm{Si})q_\mathrm{gr}]$
($q_\mathrm{Si}{=}1000$ mAh/g 원소 Si 1차 가역\cite{limthongkul2003}, tier-A$\cdot$$q_\mathrm{gr}{=}372$)로,
$m_\mathrm{Si}{=}0/10/20/30$ wt\% $\Rightarrow f_\mathrm{Si}{=}0/0.230/0.402/0.535$
($Q{=}0.97/1.26/1.62/2.09\,Q_\cell$). \textbf{(a)} 전체 창: 흑연 staging 4-peak 겹침 envelope 는
$m_\mathrm{Si}$ 증가로 정규화 높이가 $7.5\!\to\!3.7$ [1/V] 로 \emph{축소}하고($Q_\mathrm{gr}$ 절대값은
불변이나 총용량 $Q$ 가 커져 상대 몫 $1{-}f_\mathrm{Si}$ 로 줄어듦), 원소 Si 두 경사 전이(중심 $0.30/0.45$
V, 폭 $0.060/0.050$ V)의 기여가 고전위 쪽에서 \emph{성장}해 $V_n\!\approx\!0.26$ V 부근에서 순서가
역전한다. \textbf{(b)} 그 Si host 몫만 떼어 확대: $0\!\to\!1.585$
[1/V] 로 자란다($f_\mathrm{Si}$ 에 정비례). 곡선 개형(전이 중심$\cdot$폭)은 원소 Si 평균 전위대
$0.2$--$0.5$ V\cite{obrovac_chevrier2014}(tier-B) 안의 \textbf{tier-C 시연값}이며(신뢰값 아님, 피팅
override 전제), 원소 Si(elemental) 케이스는 SiO$_x$ 절대전위 같은 ``확인 필요'' 공백을 쓰지 않는다.
$f_\mathrm{Si}\!\to\!0$ 에서 (a)의 실선은 흑연 단독 평형곡선과 부동소수점까지 일치한다(bit-exact,
식~\eqref{eq:blend-limit}); \S\ref{sec:sum} 그림~\ref{fig:sumcurve} 와 동일 흑연 골격).}
\label{fig:blend-family}
\end{figure}
% 자산(v1.0.22 R7 신규): [V22-R7-F1 블렌드 f_Si wt% 스윕 가족 dQ/dV 2패널 fig:blend-family —
%   eq:blend-dqdv 실계산·elemental·총용량 정규화(면적1)·wt%→f_Si C-052·f_Si→0 bit-exact 검증분]


\subsection{정직 공백 GS-2 --- 공통-$\mu$ 완전 동시반응은 1차 근사다}\label{ssec:si-blend-gs2}
식~\eqref{eq:blend-balance} 은 \emph{평형}에서 정확하다 --- 평형이란 접촉한 모든 상이 하나의 $\mu$ 를
공유하는 상태이므로, 무한히 느린 준정적 OCV 곡선 $U_\mathrm{oc}(\bar x,T)$ 는 두 host 가 정말 공통 전위에
동시에 앉는다. 문제는 관측 dQ/dV 가 \emph{유한 율속}에서 측정된다는 데 있고, 여기서 세 실측이 완전
동시반응 가정과 정성적으로 어긋난다.

\begin{srcbox}
블렌드 실측이 보이는 편차(comp\_R4/BLEND\_UP 검증분 --- 원장 V1 등재 완료):
$\cdot$ 연속체 모델은 \emph{높은 SoC 에서 흑연이 반응전류 대부분을 담당하다가 깊은 방전으로 갈수록
급격히 Si 상으로 전환}됨을 보인다 --- 서로 다른 OCV 곡선$\cdot$질량분율$\cdot$교환전류밀도의
결과다\cite{ai_composite2022}.
$\cdot$ 고-Si 함량(30 wt\%) 블렌드를 성분 분리 측정하면 \emph{혼합 거동이 두 성분의 단순합이 아니고
(비가산적)}, 탈리튬화 시 Si$\cdot$흑연 간 유효 C-rate 격차가 특히 크다\cite{chatzogiannakis_blend2025}.
$\cdot$ SiO$_x$/흑연에서 \emph{겹치는 리튬화 전위}(overlapping lithiation potentials)가 실재해 공통-$\mu$
의 물리적 근거를 주지만, 바로 그 전위 중첩 구간의 계면에 양방향 Li$^+$ 확산이 균열을 유발한다(대가의
실측)\cite{zhan_siox2026}. 이론 다리로는 \S\ref{ssec:si-carry} 의 N9 실증(verbrugge\_lisi2016)과
같은 저자 계보(Verbrugge)의 MSMR 블렌드 축소모형이 있다\cite{tu_blend2024}.
\end{srcbox}

곧 완전 동시반응(모든 SoC 에서 두 host 가 공통 $U_\mathrm{oc}$ 에 정확히 함께 반응)은 \emph{평형 극한의
1차 근사}이고, 유한 율속에서는 (a) SoC 구간별로 우세 반응 host 가 전환되며 (b) 혼합 dQ/dV 가 두 host
평형 기여의 $f_\mathrm{Si}$-가중 단순합에서 벗어난다(비가산). 이 편차의 공백 분류는 다음과 같다.
\begin{itemize}[leftmargin=1.4em,itemsep=1pt]
\item \textbf{4분류 = 물리 가정 충돌.} 식~\eqref{eq:blend-balance} 의 ``두 host 가 한 $U_\mathrm{oc}$
에 동시반응'' 은 준정적(평형) 가정이다. 유한 율속의 host 전환$\cdot$비가산은 이 준정적 읽기와 충돌한다
--- 논리 공백(식의 결함)도, 수치해법 문제(유일근은 이미 보증됨)도, 정의상 implicit(그것은
$U_\mathrm{oc}$ 자체의 지위)도 아니다. 평형 골격은 옳게 닫혀 있고, 어긋나는 것은 그것을 유한 율속
관측에 곧바로 등치하는 \emph{물리 가정}이다.
\item \textbf{범위 밖 선언.} 구간별 host 전환$\cdot$비가산을 담으려면 host 별 교환전류밀도$\cdot$수송$\cdot$기계
구속을 넣은 반응전류 배분 모형(예: \cite{ai_composite2022} 계열)이 필요하다 --- 이는 평형 대정준 반전의
바깥, 동역학 층의 별도 하위 모형이고 본 장 범위 밖이다(마스터플랜 Non-goals: 구간별 전환 모델링은 후속
버전). §3.5 코드는 평형 합성(공통 $U_\mathrm{oc}$ 의 $f_\mathrm{Si}$-가중 합)까지만 구현하고, 유한 율속
비가산은 구현하지 않는다.
\item \textbf{경계 표식.} 따라서 본 장은 두 층을 명시 분리한다 --- \emph{평형 층}(정확: 식~\eqref{eq:blend-balance},
유일근$\cdot$$f_\mathrm{Si}\to0$ 회수 이월) 대 \emph{유한 율속 층}(1차 근사: 완전 동시반응, 편차는 GS-2
공백). 전위 중첩 구간의 실재\cite{zhan_siox2026} 가 1차 근사의 \emph{부분 지지}를, host 전환$\cdot$비가산
\cite{ai_composite2022,chatzogiannakis_blend2025} 이 그 \emph{한계}를 동시에 준다.
\end{itemize}

\begin{warnbox}
GS-2 는 골격을 부정하지 않는다 --- 식~\eqref{eq:blend-balance} 의 평형 지위와 유일근$\cdot$회수 검산은
그대로 유효하다. GS-2 가 선언하는 것은 그 평형식을 유한 율속 블렌드 dQ/dV 에 \emph{완전 동시반응으로
직결}하는 읽기가 1차 근사라는 점, 그리고 그 편차의 정량 모델링이 범위 밖이라는 점이다. 이 정직 분리가
없으면 §3.5 코드의 $f_\mathrm{Si}$-가중 합성 곡선을 실측과 1:1 로 오인할 위험이 있다.
\end{warnbox}

% 자산: [V22-R5-11 공통-μ 대정준 host 곱 인수분해 eq:blend-factor(eq:sm-mc-factor 일반화)]
%   [V22-R5-12 host×클래스 점유합 eq:blend-occ] [V22-R5-13 공통-μ 전하 보존 반전 eq:blend-balance(eq:sm-mc-balance 한 줄 일반화·중심식 유지)]
%   [V22-R5-14 요동 양성 유일근 이월+f_Si→0 bit-exact 회수 eq:blend-limit(N_p=1 회수 확장·R6 게이트 이론짝)]
%   [V22-R5-15 GS-2 비가산·host 전환 정직 공백 ssec:si-blend-gs2(물리 가정 충돌·범위 밖 선언·평형/율속 층 분리)]
